\documentclass[tikz,border=10pt]{standalone}
\usepackage{tikz}
\usetikzlibrary{positioning, arrows.meta, shapes.geometric, calc, backgrounds, fit}
\usepackage{luatexja}

\begin{document}

\begin{tikzpicture}[
    node distance=1.3cm,
    >={Latex[width=3mm,length=3mm]},
    font=\sffamily\small,
    % スタイル定義
    data block/.style={rectangle, draw=none, fill=gray!10, text centered, rounded corners, minimum height=0.8cm, minimum width=2.2cm},
    model block/.style={rectangle, draw, thick, dashed, inner sep=10pt, rounded corners},
    layer/.style={rectangle, draw=none, fill=gray!40, minimum height=0.2cm, minimum width=2.2cm},
    arrow/.style={->, thick, draw=gray!80},
    loss line/.style={<->, thick, densely dotted, draw=red!70},
    ema arrow/.style={->, very thick, draw=orange!80},
    label text/.style={font=\footnotesize, align=center}
]

% --- 色の定義 ---
\definecolor{studentColor}{RGB}{245, 183, 121} % オレンジ系
\definecolor{teacherColor}{RGB}{149, 196, 230} % 青系
\definecolor{labelColor}{RGB}{169, 209, 142}   % 緑系

% ==========================
% 1. Student Model (中央上)
% ==========================
\node[yshift=3cm] (student_center) {};

% Studentのレイヤー表現
\node[layer, fill=studentColor, minimum width=2.7cm] (s_l1) at (student_center) {};
\node[layer, fill=studentColor, minimum width=2.7cm, above=0.15cm of s_l1] (s_l2) {};
\node[layer, fill=studentColor, minimum width=2.7cm, above=0.15cm of s_l2] (s_l3) {};
% Dropoutの灰色バー (Student) - 新しい独立した棒
\node[layer, fill=gray!60, minimum width=0.15cm, minimum height=0.2cm] at ($(s_l1.west)+(0.5,0)$) {};
\node[layer, fill=gray!60, minimum width=0.15cm, minimum height=0.2cm] at ($(s_l1.west)+(1.5,0)$) {};
\node[layer, fill=gray!60, minimum width=0.15cm, minimum height=0.2cm] at ($(s_l1.west)+(2.3,0)$) {};
\node[layer, fill=gray!60, minimum width=0.15cm, minimum height=0.2cm] at ($(s_l2.west)+(0.5,0)$) {};
\node[layer, fill=gray!60, minimum width=0.15cm, minimum height=0.2cm] at ($(s_l2.west)+(1.5,0)$) {};
\node[layer, fill=gray!60, minimum width=0.15cm, minimum height=0.2cm] at ($(s_l2.west)+(2.3,0)$) {};
\node[layer, fill=gray!60, minimum width=0.15cm, minimum height=0.2cm] at ($(s_l3.west)+(0.5,0)$) {};
\node[layer, fill=gray!60, minimum width=0.15cm, minimum height=0.2cm] at ($(s_l3.west)+(1.5,0)$) {};
\node[layer, fill=gray!60, minimum width=0.15cm, minimum height=0.2cm] at ($(s_l3.west)+(2.3,0)$) {};

% Student枠
\node[model block, draw=studentColor!80, fit=(s_l1)(s_l3), label={[text=studentColor!80]below:\textbf{Student Model}}] (student_box) {};

% ==========================
% 2. Teacher Model (右側上)
% ==========================
\node[right=5.5cm of student_center] (teacher_center) {};

% Teacherのレイヤー表現
\node[layer, fill=teacherColor, minimum width=2.7cm] (t_l1) at (teacher_center) {};
\node[layer, fill=teacherColor, minimum width=2.7cm, above=0.15cm of t_l1] (t_l2) {};
\node[layer, fill=teacherColor, minimum width=2.7cm, above=0.15cm of t_l2] (t_l3) {};
% Dropoutの灰色バー (Teacher) - Studentと同じ位置
\node[layer, fill=gray!60, minimum width=0.15cm, minimum height=0.2cm] at ($(t_l1.west)+(0.5,0)$) {};
\node[layer, fill=gray!60, minimum width=0.15cm, minimum height=0.2cm] at ($(t_l1.west)+(1.5,0)$) {};
\node[layer, fill=gray!60, minimum width=0.15cm, minimum height=0.2cm] at ($(t_l1.west)+(2.3,0)$) {};
\node[layer, fill=gray!60, minimum width=0.15cm, minimum height=0.2cm] at ($(t_l2.west)+(0.5,0)$) {};
\node[layer, fill=gray!60, minimum width=0.15cm, minimum height=0.2cm] at ($(t_l2.west)+(1.5,0)$) {};
\node[layer, fill=gray!60, minimum width=0.15cm, minimum height=0.2cm] at ($(t_l2.west)+(2.3,0)$) {};
\node[layer, fill=gray!60, minimum width=0.15cm, minimum height=0.2cm] at ($(t_l3.west)+(0.5,0)$) {};
\node[layer, fill=gray!60, minimum width=0.15cm, minimum height=0.2cm] at ($(t_l3.west)+(1.5,0)$) {};
\node[layer, fill=gray!60, minimum width=0.15cm, minimum height=0.2cm] at ($(t_l3.west)+(2.3,0)$) {};

% Teacher枠
\node[model block, draw=teacherColor!80, fit=(t_l1)(t_l3), label={[text=teacherColor!80, name=teacher_label]below:\textbf{Teacher Model}}] (teacher_box) {};

% ==========================
% 3. 入力データ (下側)
% ==========================
\node[data block, fill=labelColor, below=1.2cm of student_box] (weak_label) {ラベル付き音声};
\node[data block, fill=yellow!20, right=0.8cm of weak_label] (unlabeled) {ラベルなし音声};

% ==========================
% 4. 出力: フレームレベル (Frame-level)
% ==========================
% Student Frame Output
\node[above=0.6cm of student_box] (s_frame_base) {};
\foreach \x in {-1.2, -0.9, ..., 1.2} {
    \draw[fill=studentColor, draw=none] ($(s_frame_base)+(\x,0)$) rectangle ++(0.25, {rand*0.3+0.4});
}
\draw[gray] ($(s_frame_base)+(-1.3,0)$) -- ($(s_frame_base)+(1.55,0)$);
\node[label text, left=1.4cm of s_frame_base] (s_frame_lbl) {強予測 確率};

% Teacher Frame Output
\node[above=0.6cm of teacher_box] (t_frame_base) {};
\foreach \x in {-1.2, -0.9, ..., 1.2} {
    \draw[fill=teacherColor, draw=none] ($(t_frame_base)+(\x,0)$) rectangle ++(0.25, {rand*0.3+0.4});
}
\draw[gray] ($(t_frame_base)+(-1.3,0)$) -- ($(t_frame_base)+(1.55,0)$);

% ==========================
% 5. プーリング (Pooling)
% ==========================
\draw[->, thick, gray!60] ($(s_frame_base)+(0,0.6)$) -- ($(s_frame_base)+(0,1.2)$);

\draw[->, thick, gray!60] ($(t_frame_base)+(0,0.6)$) -- ($(t_frame_base)+(0,1.2)$);

% ==========================
% 6. 出力: クリップレベル (Clip-level)
% ==========================
% Student Clip Output (Bar chart)
\node[above=1.2cm of s_frame_base] (s_clip_base) {};
\draw[fill=studentColor, draw=none] ($(s_clip_base)+(-1,0)$) rectangle ++(0.3, 0.8);
\draw[fill=studentColor, draw=none] ($(s_clip_base)+(-0.5,0)$) rectangle ++(0.3, 0.2);
\draw[fill=studentColor, draw=none] ($(s_clip_base)+(0,0)$) rectangle ++(0.3, 0.9);
\draw[fill=studentColor, draw=none] ($(s_clip_base)+(0.5,0)$) rectangle ++(0.3, 0.4);
\draw[gray] ($(s_clip_base)+(-1.2,0)$) -- ($(s_clip_base)+(1.2,0)$);
\node[label text, left=1.4cm of s_clip_base] (s_clip_lbl) {弱予測 確率};

% Teacher Clip Output
\node[above=1.2cm of t_frame_base] (t_clip_base) {};
\draw[fill=teacherColor, draw=none] ($(t_clip_base)+(-1,0)$) rectangle ++(0.3, 0.8);
\draw[fill=teacherColor, draw=none] ($(t_clip_base)+(-0.5,0)$) rectangle ++(0.3, 0.25);
\draw[fill=teacherColor, draw=none] ($(t_clip_base)+(0,0)$) rectangle ++(0.3, 0.85);
\draw[fill=teacherColor, draw=none] ($(t_clip_base)+(0.5,0)$) rectangle ++(0.3, 0.35);
\draw[gray] ($(t_clip_base)+(-1.2,0)$) -- ($(t_clip_base)+(1.2,0)$);

% ==========================
% 7. 正解ラベル (Ground Truth)
% ==========================
\node[above=1.0cm of s_clip_base] (gt_base) {};
\node[data block, fill=labelColor!30, minimum width=2.2cm, minimum height=0.6cm] at (gt_base) (gt_box) {正解ラベル};

% ==========================
% 7.5. 弱予測と強予測を囲む枠
% ==========================
\node[draw=red!50, densely dotted, thick, rounded corners, fit=(s_clip_lbl)(s_frame_lbl), inner sep=5pt] (pred_group) {};

% ==========================
% 8. 矢印と接続
% ==========================

%Input -> Models
\draw[arrow] (weak_label.north) -- ($(student_box.south)+(0,-0.3)$);
% ラベルなし音声から両モデルへの分岐（線を結合）
\coordinate (split_point) at ($(unlabeled.north)+(0,0.3)$);
\draw[thick, draw=gray!80] (unlabeled.north) -- (split_point) -| ($(student_box.south)+(0.5,-0.3)$);
\draw[thick, draw=gray!80] (split_point) -| ($(teacher_box.south)+(0,-0.3)$);
% 最終接続点に矢印を追加
\draw[arrow] ($($(student_box.south)+(0.5,-0.3)$)-(0,0.1)$) -- ($(student_box.south)+(0.5,-0.3)$);
\draw[arrow] ($($(teacher_box.south)+(0,-0.3)$)-(0,0.1)$) -- ($(teacher_box.south)+(0,-0.3)$);

% Models -> Outputs
\draw[arrow] (student_box.north) -- (s_frame_base.south);
\draw[arrow] (teacher_box.north) -- (t_frame_base.south);

% EMA Update
\draw[ema arrow] (student_box.east) -- node[above, font=\bfseries\footnotesize, align=center] {EMA更新} (teacher_box.west);

% ==========================
% 9. ロス（コスト）の表示
% ==========================

% Classification Cost
\draw[loss line] (s_clip_lbl.north) to[bend left=20] node[left, font=\bfseries, text=red!70] {通常の損失} (gt_box.west);

% Consistency Cost (中央に1本: 弱予測と強予測の間)
\coordinate (consistency_left) at ($(s_clip_base)!0.5!(s_frame_base) + (1.5,0)$);
\coordinate (consistency_right) at ($(t_clip_base)!0.5!(t_frame_base) + (-1.5,0)$);
\draw[loss line] (consistency_left) -- node[above, font=\bfseries, text=red!70, fill=white] {一貫性損失} (consistency_right);

\end{tikzpicture}
\end{document}